\documentclass[a4paper,landspace]{article}
\usepackage[paperwidth=14.13in, paperheight=9.86in]{geometry}
%% \usepackage{geometry}

\usepackage[usenames,dvipsnames]{xcolor}
\usepackage[utf8]{inputenc}
\usepackage{fontspec}
\usepackage{tikz}
\usetikzlibrary{positioning}
\usepackage{rotating}
\usepackage[absolute,overlay]{textpos}
%% Useful to print in an A4/A3 sheet and see the actual size of paper
%% \usepackage[frame,a3,landscape,pdftex,center]{crop}

% donot indent first line of paragraph.
\setlength{\parindent}{0pt}
\setlength{\parskip}{5pt plus 2pt minus 1pt}

\sloppy

%% Completely remove hyphenation
\usepackage[none]{hyphenat}
\usepackage[most]{tcolorbox}
\usepackage{setspace}

\begin{document}
\newfontfamily{\cabinbold}
              [ Path = ../fonts/,
                UprightFont = Cabin-Bold.ttf,
              ]{ Cabin }
\newfontfamily{\dejavusans}
              [ Path = ../fonts/,
                UprightFont = DejaVuSans.ttf,
                BoldFont = DejaVuSans-Bold.ttf,
              ]{ DejaVu Sans }
\newgeometry{top=0cm, bottom=0cm}
\pagestyle{empty}
\pagecolor{Blue}

\begin{textblock*}{0in}(6.97in,0.5in) % {block width} (coords)
  \begin{turn}{-90}
    {\fontsize{16}{16} {\cabinbold {\color{Yellow} Essentials of Go Programming}}}
  \end{turn}
\end{textblock*}

\begin{textblock*}{0in}(6.99in,5.0in) % {block width} (coords)
  \begin{turn}{-90}
    {\fontsize{14}{14} {\cabinbold {\color{White} First Edition}}}
  \end{turn}
\end{textblock*}

\begin{textblock*}{0in}(6.97in,7.5in) % {block width} (coords)
  \begin{turn}{-90}
    {\fontsize{16}{16} {\cabinbold {\color{White} Baiju Muthukadan}}}
  \end{turn}
\end{textblock*}

\begin{textblock*}{5in}(8in,1.7in) % {block width} (coords)
    {\fontsize{50}{50} {\cabinbold {\color{Yellow} Essentials of}}}
\end{textblock*}

\begin{textblock*}{5in}(8in,2.5in) % {block width} (coords)
    {\fontsize{58}{58} {\cabinbold {\color{Yellow} Go Programming}}}
\end{textblock*}

\begin{textblock*}{5in}(10.8in,3.59in) % {block width} (coords)
    {\fontsize{36}{36} {\cabinbold {\color{White} First Edition}}}
\end{textblock*}

\begin{textblock*}{5in}(8.25in,4.4in) % {block width} (coords)

%% Credit: https://commons.wikimedia.org/wiki/File:The_ladder_of_life_is_full_of_splinters.jpg
\tcbox{\includegraphics[width=4.7in,keepaspectratio]{images/ladder.jpg}}

\end{textblock*}

\begin{textblock*}{5in}(8.59in,8.6in) % {block width} (coords)
    {\fontsize{42}{42} {\cabinbold {\color{White} Baiju Muthukadan}}}
\end{textblock*}

\begin{textblock*}{5.6in}(0.5in,0.4in) % {block width} (coords)
  {\fontsize{12}{12}

    {\dejavusans {\setstretch{1.2}\color{White}Having proficiency in computer
        programming empowers you to tackle a wide range of real-world
        challenges. Programming encompasses a process that begins with
        formulating a computing problem and culminates in the creation of
        computer programs or software. It is an integral part of a broader
        software development process.

Programming entails analyzing, designing, and implementing software solutions.
Coding is the active process of putting the software into practice. At times,
coding involves utilizing multiple programming languages and incorporating other
technologies. Acquiring proficiency in a programming language is an essential
aspect of the skill set required for computer programming and software
development as a whole.

By utilizing a programming language, we construct instructions for computing
machines such as desktop computers, laptops, and mobile phones.

Today, there exists a multitude of programming languages with diverse features.
Choosing the appropriate programming language for a given problem is crucial.
Certain languages are better suited for specific challenges. This book serves as
an introduction to the Go programming language, offering valuable insights to
aid in making informed language selection decisions for your projects. If you
have already chosen to work with Go and seek to enhance your understanding, this
book provides an excellent introduction to Go programming.

Go, an open-source language developed by Google with significant contributions
from the community, has emerged as a popular choice. The project was initiated
in 2007 by Robert Griesemer, Rob Pike, and Ken Thompson and was subsequently
released as open-source software by Google in November 2009. Go has gained
widespread adoption across diverse problem domains, being embraced by numerous
organizations. }}}
\end{textblock*}

\begin{textblock*}{5in}(1.2in,5.8in) % {block width} (coords)
\begin{tikzpicture}
\draw [Yellow, line width=0.1in, dashed] (0,0) -- (10.3,0);
\end{tikzpicture}
\end{textblock*}

\begin{textblock*}{5.6in}(0.5in,6.0in) % {block width} (coords)
  {\fontsize{12}{12}

    {\dejavusans {\setstretch{1.2}\color{White}Baiju Muthukadan is a software
        engineer from Kerala, India. He began his software development career in
        2002. Over the years, Baiju has contributed to many free and open source
        software projects. In 2007, he authored his first book on Python. Since
        2013, he has been working on various Go projects. His Twitter handle is
        @baijum\\ }}

    {\dejavusans {\color{White}
        Source repository: https://github.com/baijum/essential-go}} }

\end{textblock*}

\end{document}
